
Permutation equivariant networks for weight space learning, such as Neural Functional Networks (NFN) \citep{Zhouetal2023}, exploit the fact that any permutation of neurons in a hidden layer, paired with the inverse permutation on the adjacent layer, leaves the network function unchanged. This approach yields Kendall's tau > 0.93 on CNN generalization benchmarks \citep{Zhouetal2023}. However, the same permutation equivariant methods perform poorly when applied zero-shot to transformer architectures. Preliminary evaluation in this work indicates tau falls below 0.50 in such cross-architecture settings, representing a gap of approximately 43 percentage points relative to within-architecture CNN performance. The cause of this gap is not definitively established in prior work.

Two main approaches exist for cross-architecture WSL. Architecture-specific equivariant networks (NFN, Transformer-NFN) excel within a given architecture family but are not designed for zero-shot cross-architecture transfer [Zhou et al., 2023; Tran-Viet et al., 2024]. Token-based representations such as SANE \citep{Schurholtetal2024} operate across architectures by treating weight rows as tokens with sequential positional encodings, but do not encode symmetry structure. Neither approach has characterized which symmetry groups explain the dominant weight-space variation for each layer type.

This paper addresses that measurement gap. The input/output channel permutation group G = S_{c_in} x S_{c_out} acts on all standard linear operator types (Conv2d, Linear, MultiheadAttention) as a consequence of their shared linear operator structure. The present work verifies this symmetry empirically and measures what fraction of weight-space variance it captures, stratified by layer type, across a large model zoo.

The research is organized around three sub-hypotheses in a sequential gate-controlled design (H-E1, H-M1, H-M2), where each downstream experiment requires the preceding gate to pass. A fourth sub-hypothesis (H-M3, cross-architecture training) and a fifth (H-C1, OVR decomposition) were pre-specified but were not executed because H-M2 failed its gate criterion.

The three executed sub-hypotheses yield the following findings:

\begin{enumerate}
\item \textbf{Exact symmetry validation (H-E1):} |Delta acc| = 0.000000 across 4,500 permutation runs on CNN Zoo and Transformer Zoo checkpoints, satisfying the gate threshold by several orders of magnitude.
\end{enumerate}

\begin{enumerate}
\item \textbf{Practical computability (H-M1):} Orbit-PE is computable for all layer types with 1.167x mean overhead (threshold: <= 1.2x), using a unified codebase (HAS_ARCH_BRANCHES = False) and consistent output dimensionality (token_dim = 64).
\end{enumerate}

\begin{enumerate}
\item \textbf{Zoo-scale variance stratification (H-M2):} The overall permutation-orbit variance fraction is 0.3479 +/- 0.0536 (n = 1,000 models x 50 epochs), below the 0.60 gate threshold. Layer-type breakdown reveals Conv2d ratio = 0.637 and Linear (FC) ratio = 0.133 — a 4.8x difference. The gate FAILED, blocking H-M3 execution. A pre-specified pivot to hybrid orbit-PE encoding was activated.
\end{enumerate}

The paper is organized as follows. Section 2 surveys related work. Section 3 describes the methodology. Section 4 details the experimental setup. Section 5 presents results. Section 6 discusses implications and limitations. Section 7 concludes.

